\section{Introduction}
\label{sec:intro}

Imagine a user who moves from Shanghai to Beijing and then asks their AI assistant for a dinner recommendation. A well-designed assistant should not suggest their old favorite restaurant in Shanghai. Yet current memory systems — including MemGPT~\cite{packer2024memgpt}, Mem0, and PersonaMem~\cite{jiang2025personamem} — store facts as flat key-value pairs and update memories only when the user explicitly says ``I no longer like that.'' Without explicit retraction, stale facts persist indefinitely, and AI assistants continue making recommendations based on outdated user states.

This failure is not limited to explicit location changes. Consider a user who has gradually drifted toward quiet evenings at home over the past few months. No single conversation contains an explicit statement like ``I no longer enjoy bars.'' Instead, the drift is revealed through repeated mentions of ``quiet evenings,'' ``reading at home,'' and ``early nights.'' An assistant that stores only stated facts will miss this drift entirely, continuing to recommend bars and social events the user has outgrown.

{\bf What makes memory invalidation hard?} Two challenges combine to make this difficult. First, a \textit{structural challenge}: when the root context changes (the user moves cities), the system must identify all memories that depended on that context. Second, a \textit{semantic challenge}: the system must establish a conflict relationship between ``prefers quiet'' and ``likes bars'' without the user ever making this connection explicit. We call the second challenge the \textit{semantic bridging problem}.

\para{Existing work falls short on both counts.} \horizonbench~\cite{li2026horizonbench} demonstrates that all 25 evaluated frontier models suffer from belief-update failure, selecting stale pre-evolution preferences as answers more than 25\% of the time. LoCoMo~\cite{maharana2024locomo} provides real-world conversational data for long-term memory evaluation but lacks cascade invalidation methodology. PAMU~\cite{sun2025pamu} addresses preference drift detection on LoCoMo using sliding-window statistics but does not propagate invalidation through a graph. RippleEdits~\cite{cohen2023rippleedits} evaluates knowledge editing cascades but focuses on factual encyclopedic knowledge, not user preferences. No prior work evaluates cascading invalidation on \textit{real} conversational data using both structural and semantic conflict edges.

\para{We propose cascading memory invalidation via typed dependency edges.} We model each user's memories as nodes in a directed graph $\gG = (\gV, \gE)$ with two typed edge classes: \locatedin edges (from a root location node to location-dependent memory nodes) and directed \conflictswith edges (from a current preference node to an outdated one). When a root context changes, a weighted BFS propagation automatically downweights all transitively dependent memories, while leaving unrelated memories untouched.

The core novel insight is edge \textit{directionality}: \conflictswith edges must flow \textbf{from current to old}, not bidirectionally. We empirically demonstrate that bidirectional edges cause mutual cancellation, reducing full-cascade performance to the flat-memory baseline (41.8\% vs.\ 72.2\% with directed edges). This failure mode has not been documented in prior work.

\para{Our results show that all four research hypotheses are supported.} We evaluate on \locomo (35 real multi-session dialogues) and \horizonbench (4,245 preference evolution items). Structural cascade via \locatedin edges achieves 81.7\% accuracy on location-shift events. LLM-based \conflictswith edge detection achieves 100\% precision and 90.9\% recall. Full transitive cascade improves evolved-preference accuracy by +32.2 percentage points over flat memory (72.2\% vs.\ 40.1\%), and reduces stale-preference selection from 43.2\% to 15.4\%.

Our contributions are:

\begin{itemize}[leftmargin=*,itemsep=0pt,topsep=0pt]
    \item We propose a directed memory graph with typed \locatedin and \conflictswith edges that enable automated cascading invalidation without explicit user retraction.
    \item We demonstrate that \gptfouromini solves the semantic bridging problem with 100\% precision at memory write time, establishing \conflictswith edges from world knowledge about preference conflicts.
    \item We identify and empirically characterize \textit{bidirectional edge cancellation} as a critical failure mode in graph-based memory invalidation, providing a concrete design principle for future systems.
    \item We conduct comprehensive evaluations on \locomo and \horizonbench, showing +32.2pp improvement over flat memory and 64\% relative reduction in stale-preference selection rate.
\end{itemize}
